\section{External CPTAC mutation-inference audit}

To test cohort transport without external-label tuning, a fixed logistic baseline was fitted on TCGA-PAAD pathway and mutation-burden features, excluding the target driver. The same feature definitions were then applied to the CPTAC-PDAC mutation profile (140 samples, 4,897 variant records). This is not external validation of the previous GNN/CNN. cBioPortal sample identifiers do not overlap, but distinct study identifiers alone do not establish patient identity independence.

\begin{center}\begin{tabular}{lrrr}\hline Target & Positive / 140 & AUROC & Bootstrap 95\%\\\hline

KRAS & 135 & 0.384 & 0.157--0.621 \\

TP53 & 105 & 0.741 & 0.643--0.830 \\

CDKN2A & 29 & 0.497 & 0.376--0.612 \\

SMAD4 & 24 & 0.347 & 0.253--0.449 \\

\hline\end{tabular}\end{center}

TP53 provides a positive transport signal, but the CDKN2A endpoint is at chance and SMAD4 reverses direction. KRAS has only five mutation-negative test samples, so its uncertainty is large. These failures prevent claiming cohort-general CDKN2A interpretation. Mutation-only targets omit deletion and methylation, and purity/variant-calling differences remain unmodeled. Confidence intervals are stratified bootstraps of fixed model predictions; they do not include refitting uncertainty or cohort uncertainty.

The experiment supports a narrower contribution: a reproducible external audit that distinguishes a transferable TP53 association from nontransferable driver signals. It establishes neither a mechanism nor clinical utility. Source: \url{https://www.cbioportal.org/study/summary?id=paad_cptac_2021}. Input hashes, prevalence baselines and calibration errors are preserved in results/cptac\_external\_validation.json.
\subsection{External discrimination does not generally beat a training-prior probability constant}
A frozen proper-score unit refits the identical TCGA-only balanced logistic baseline and verifies the earlier CPTAC AUROC, average precision and Brier values to $10^{-10}$. No CPTAC label is used for fitting, shift selection or calibration. The comparator constant is each target's TCGA training prevalence, not its CPTAC prevalence. A second comparator applies training-prior odds shifting to the balanced prediction. This preserves all AUROC and average precision values, but is a heuristic rather than certified probability calibration.

\begin{center}
\begin{tabular}{lrrr}
\hline
Target & Balanced Brier & Train-prior constant & Train-prior shift \\
\hline
KRAS & 0.255 & 0.142 & 0.186 \\
TP53 & 0.185 & 0.216 & 0.170 \\
CDKN2A & 0.325 & 0.165 & 0.182 \\
SMAD4 & 0.318 & 0.143 & 0.165 \\
\hline
\end{tabular}
\end{center}
Only TP53 beats the training-prior constant, on both Brier score and log loss. Its log losses are 0.557 balanced, 0.624 constant and 0.518 shifted. KRAS has 0.778, 0.473 and 0.596; CDKN2A has 0.903, 0.511 and 0.648; SMAD4 has 0.833, 0.461 and 0.517. The odds shift improves both proper scores for all four targets relative to balanced outputs, yet three still lose to the simple source-prior constant. Thus the TP53 discrimination result survives this probability-score comparator, while the other targets do not acquire useful probability evidence merely from an improved intercept heuristic.

TCGA and CPTAC prevalences differ: KRAS 0.636 versus 0.964, TP53 0.582 versus 0.750, CDKN2A 0.190 versus 0.207 and SMAD4 0.201 versus 0.171. Those target-cohort values are descriptive, never used to alter predictions. KRAS has only five negatives, so neither this finding nor its poor AUROC is a reliable subgroup clinical estimate. No patient-identity independence, variant-platform harmonization, deletion or methylation target, GNN validation, diagnosis, treatment or clinical utility is established.

All 560 target/sample predictions, fixed source priors and proper scores are in \path{results/cptac_probability_audit.json}, with source hashes and protocol in \path{scripts/cptac_probability_plan.json}. Source: \url{https://www.cbioportal.org/study/summary?id=paad_cptac_2021}. The content-page gate remains uncertified.

